% lagrangian-product-algorithm-proof.tex — \input'd from proofs.tex
% Conventions: .claude/rules/thesis-tex.md
%
% Sources: Bezdek-Bezdek (2009), Rudolf (2022), Artstein-Avidan & Ostrover (2014),
%   HK2017. Agent-written, math verified by Joern.
%
% Structure: Lagrangian products, orbit structure, billiard characterization

% Jörn: math approved (e1345bb) — from \subsection{Lagrangian Products} to end of subsection

\begin{definition}[Lagrangian subspaces]\label{def:lagrangian-subspaces}
% QC: L_q, L_p are Lagrangian because omega_0 vanishes on each.
%   dim L_q = dim L_p = 2, and R^4 = L_q + L_p (direct sum).
% Downstream: L_q = components [0,1], L_p = components [2,3].
%   is_q_type(n) = |n[2]| < eps && |n[3]| < eps.
%   is_p_type(n) = |n[0]| < eps && |n[1]| < eps.
The \emph{$q$-subspace} $L_q \coloneqq \{(q, 0) : q \in \mathbb{R}^2\}$
and \emph{$p$-subspace} $L_p \coloneqq \{(0, p) : p \in \mathbb{R}^2\}$
are Lagrangian subspaces of $(\mathbb{R}^4, \omega_0)$:
\[
  \omega_0|_{L_q} = 0, \qquad
  \omega_0|_{L_p} = 0, \qquad
  \mathbb{R}^4 = L_q \oplus L_p.
\]
\end{definition}

\begin{definition}[Lagrangian product]\label{def:lagrangian-product}
% QC: K_q subset L_q and K_p subset L_p, both 2D convex bodies
%   with 0 in their interior.
%   K = K_q x K_p is a 4D polytope with 0 in int(K).
%   Facet count: F = F_q + F_p where F_q, F_p are edge counts
%   of K_q, K_p respectively.
% Downstream: construct by embedding 2D normals:
%   q-facet (a,b) -> (a, b, 0, 0), p-facet (c,d) -> (0, 0, c, d).
%   Heights preserved. Vertices: all (v_q, v_p) for v_q in vert(K_q), v_p in vert(K_p).
Let $K_q \subset L_q$ and $K_p \subset L_p$ be convex polygons
with $0 \in \operatorname{int}(K_q)$ and $0 \in \operatorname{int}(K_p)$.
The \emph{Lagrangian product} is the polytope
\[
  K \coloneqq K_q \times K_p
  = \{ (q, p) \in \mathbb{R}^4 : q \in K_q,\; p \in K_p \}.
\]
\end{definition}

\begin{lemma}[Facet structure of Lagrangian products]\label{lem:lagrangian-facets}
% QC: every facet is either q-type or p-type, never mixed.
%   q-type: normal (n_q, 0), height h_q = h_{K_q}(n_q).
%   p-type: normal (0, n_p), height h_p = h_{K_p}(n_p).
%   Facet count F = F_q + F_p (edges of K_q + edges of K_p).
% Downstream: classify_facets checks |n[2..3]| < eps or |n[0..1]| < eps.
Let $K = K_q \times K_p$ be a Lagrangian product.
Let $K_q$ have edges with outward unit normals
$\hat{n}_1^q, \ldots, \hat{n}_{F_q}^q \in S^1 \subset L_q$
and heights $h_1^q, \ldots, h_{F_q}^q > 0$,
and similarly $\hat{n}_1^p, \ldots, \hat{n}_{F_p}^p$ and $h_j^p$ for~$K_p$.
Then $K$ has $F = F_q + F_p$ facets, partitioned into:
\begin{itemize}
\item \emph{$q$-facets} with normals $n_i = (\hat{n}_i^q, 0)$
  and heights $h_i = h_i^q$, for $i = 1, \ldots, F_q$;
\item \emph{$p$-facets} with normals $n_j = (0, \hat{n}_j^p)$
  and heights $h_j = h_j^p$, for $j = 1, \ldots, F_p$.
\end{itemize}
Every facet normal lies entirely in $L_q$ or entirely in $L_p$.
\end{lemma}

\begin{proof}
A facet of $K_q \times K_p$ arises from a facet (edge) of
one factor and the full other factor: either
$\text{edge}_i(K_q) \times K_p$ or $K_q \times \text{edge}_j(K_p)$.
The outward normal to $\text{edge}_i(K_q) \times K_p$ is
$(\hat{n}_i^q, 0)$ with height $h_{K_q}(\hat{n}_i^q) = h_i^q$;
similarly for $p$-facets.
\end{proof}

\begin{lemma}[Reeb vectors on Lagrangian products]\label{lem:lagrangian-reeb}
% QC: R_i = (2/h_i) J_0 n_i per def:reeb-vectors.
%   q-facet: n = (n_q, 0), so J_0 n = (0, n_q), so R = (2/h)(0, n_q) in L_p.
%   p-facet: n = (0, n_p), so J_0 n = (-n_p, 0), so R = (2/h)(-n_p, 0) in L_q.
% Downstream: on q-facets, only p-coordinates change.
%   On p-facets, only q-coordinates change.
On a Lagrangian product $K = K_q \times K_p$:
\begin{enumerate}
\item For a $q$-facet with normal $n = (\hat{n}^q, 0)$ and height~$h$:
  \[
    R = \frac{2}{h}\, J_0\, (\hat{n}^q, 0) = \frac{2}{h}\, (0, \hat{n}^q).
  \]
  The Reeb vector lies in $L_p$: the orbit moves in the $p$-direction,
  and $q$ is constant.

\item For a $p$-facet with normal $n = (0, \hat{n}^p)$ and height~$h$:
  \[
    R = \frac{2}{h}\, J_0\, (0, \hat{n}^p) = \frac{2}{h}\, (-\hat{n}^p, 0).
  \]
  The Reeb vector lies in $L_q$: the orbit moves in the $q$-direction,
  and $p$ is constant.
\end{enumerate}
\end{lemma}

\begin{proof}
Apply Definition~\ref{def:reeb-vectors} with
$J_0(q, p) = (-p, q)$
(Definition~\ref{def:symplectic-form}).
\end{proof}

\begin{lemma}[Symplectic form vanishes on same-type pairs]\label{lem:omega-vanishes}
% QC: omega_0(q, q) = 0 and omega_0(p, p) = 0 because L_q and L_p are Lagrangian.
%   omega_0(q, p) = n_q1 * n_p1 + n_q2 * n_p2 (inner product of 2D normals)
%   in general nonzero.
% Downstream: in the H matrix, same-type entries are zero.
%   Q(beta) only gets contributions from cross-type pairs.
For normals of a Lagrangian product:
\begin{enumerate}
\item $\omega_0\bigl((\hat{n}^q, 0),\; (\hat{m}^q, 0)\bigr) = 0$
  for any two $q$-facet normals.
\item $\omega_0\bigl((0, \hat{n}^p),\; (0, \hat{m}^p)\bigr) = 0$
  for any two $p$-facet normals.
\item $\omega_0\bigl((\hat{n}^q, 0),\; (0, \hat{n}^p)\bigr)
  = \langle \hat{n}^q, \hat{n}^p \rangle$,
  which is generically nonzero.
\end{enumerate}
\end{lemma}

\begin{proof}
From Definition~\ref{def:symplectic-form}:
$\omega_0(u, v) = u_{q_1} v_{p_1} - u_{p_1} v_{q_1} + u_{q_2} v_{p_2} - u_{p_2} v_{q_2}$.
If both vectors lie in $L_q$ (zero $p$-components), every term vanishes.
Similarly for $L_p$.
For $u = (\hat{n}^q, 0)$ and $v = (0, \hat{n}^p)$:
$\omega_0(u, v) = \hat{n}^q_1 \hat{n}^p_1 + \hat{n}^q_2 \hat{n}^p_2
= \langle \hat{n}^q, \hat{n}^p \rangle$.
\end{proof}


\begin{lemma}[Adjacency in Lagrangian products]\label{lem:lagrangian-adjacency}
% QC: vertices of K_q x K_p are exactly (v_q, v_p) for v_q in vert(K_q),
%   v_p in vert(K_p). Every q-facet contains all (v_q, *) where v_q is on
%   the q-edge, so it shares vertices with every p-facet.
% Downstream: in the adjacency matrix, adj[q-facet][p-facet] = true always.
%   adj[q_i][q_j] = true iff q-edges i,j share a vertex in K_q.
%   adj[p_i][p_j] = true iff p-edges i,j share a vertex in K_p.
In the facet adjacency graph $G_K$
(Definition~\ref{def:adjacency-graph}) of a Lagrangian product:
\begin{enumerate}
\item Every $q$-facet is adjacent to every $p$-facet.
\item Two $q$-facets are adjacent if and only if
  the corresponding edges of $K_q$ share a vertex.
\item Two $p$-facets are adjacent if and only if
  the corresponding edges of $K_p$ share a vertex.
\end{enumerate}
\end{lemma}

\begin{proof}
The vertices of $K_q \times K_p$ are exactly the pairs
$(v_q, v_p)$ with $v_q \in \operatorname{vert}(K_q)$,
$v_p \in \operatorname{vert}(K_p)$.
The $q$-facet with normal $(\hat{n}_i^q, 0)$ contains
$(v_q, v_p)$ whenever $v_q$ lies on edge~$i$ of~$K_q$;
since edge~$i$ has at least two vertices, and $K_p$ has at
least three vertices, every $q$-facet contains at least six
4D vertices. Every $p$-facet likewise.
Two facets are adjacent iff they share a vertex
(Definition~\ref{def:adjacency-graph}).

(i) Any $q$-facet (edge~$i$ of $K_q$) contains $(v_q, v_p)$
for some vertex $v_q$ on edge~$i$. Any $p$-facet (edge~$j$
of $K_p$) contains $(v_q', v_p)$ for some vertex $v_p$ on
edge~$j$. Since both factors range over all partner vertices,
$(v_q, v_p)$ lies on both facets for any $v_q$ on edge~$i$
and $v_p$ on edge~$j$. Adjacent.

(ii)--(iii) Two $q$-facets share a vertex $(v_q, v_p)$ iff
$v_q$ lies on both $q$-edges, i.e.\ the edges share the
vertex $v_q$ in~$K_q$. Symmetric for $p$.
\end{proof}

\begin{lemma}[Sigma structure for Lagrangian products]\label{lem:sigma-structure}
% QC: this is the central structural result. It restricts the
%   space of (S, sigma) to test for Lagrangian products.
% Downstream: enumerate sigma matching ([Q|QQ][P|PP])^k only.
%   Block size at most 2. Adjacent facets within each block.
Let $K = K_q \times K_p$ be a Lagrangian product, and let
$(\sigma, \tau)$ be a simple Reeb orbit on $\partial K$
(Definition~\ref{def:simple-reeb-orbit}) with all $\tau_k > 0$.
Write the cyclic sequence $\sigma$ as a word in the alphabet
$\{Q, P\}$ by replacing each $q$-facet index with $Q$ and each
$p$-facet index with $P$.
Then:
\begin{enumerate}
\item[(i)] The word has the form
  $\bigl([Q \mid QQ]\,[P \mid PP]\bigr)^k$ for some $k \geq 1$:
  alternating blocks of $q$- and $p$-facets, each block of length
  one or two.

\item[(ii)] Within each block of length two, the two facets are
  adjacent in $G_K$
  (consecutive edges of $K_q$ or $K_p$ sharing a vertex).
\end{enumerate}
\end{lemma}

\begin{proof}
\textbf{Claim: no three consecutive same-type facets.}\quad
Suppose for contradiction that $\sigma$ contains three
consecutive $q$-facets $q_1, q_2, q_3$ (all with $\tau > 0$).
By Lemma~\ref{lem:lagrangian-reeb}(i), the Reeb vectors
$R_{q_1}, R_{q_2}, R_{q_3}$ all lie in~$L_p$, so $q$ is
constant throughout the three segments.
By Lemma~\ref{lem:adjacency-constraint}, the orbit at the
breakpoint between $q_1$ and $q_2$ lies on both facets,
meaning $\langle q, \hat{n}_1^q \rangle = h_1^q$ and
$\langle q, \hat{n}_2^q \rangle = h_2^q$: the point $q$
lies on both edges $1$ and $2$ of~$K_q$, hence at their
common vertex $v_{12}$.
Similarly, the breakpoint between $q_2$ and $q_3$ forces $q$
to lie at vertex $v_{23}$ (intersection of edges $2$ and $3$).
Since $q$ is constant, $v_{12} = v_{23}$.
But for a nondegenerate convex polygon, $v_{12}$ and $v_{23}$
are the two distinct endpoints of edge~$2$: contradiction.

The same argument applies to $p$-facets by symmetry
(Lemma~\ref{lem:lagrangian-reeb}(ii), with $p$ constant).

\medskip\noindent
\textbf{Alternation and block structure.}\quad
Since no three consecutive facets have the same type, each maximal
run of same-type facets has length one or two.
Consecutive runs must alternate between $q$ and $p$ (a $q$-run
followed by another $q$-run would create three or more consecutive
$q$-facets at their boundary, contradicting the above).
% QC: the cyclic structure ensures equal counts: if there are
%   k q-blocks and l p-blocks with k != l, the alternation
%   would fail at the cyclic wraparound.
In a cyclic sequence, alternation forces equal block counts:
$k$ blocks of $q$-facets and $k$ blocks of $p$-facets.

\medskip\noindent
\textbf{Within-block adjacency.}\quad
Follows directly from Lemma~\ref{lem:adjacency-constraint}:
consecutive facets in $\sigma$ share a breakpoint, hence are
adjacent in~$G_K$.
By Lemma~\ref{lem:lagrangian-adjacency}(ii)--(iii),
two same-type adjacent facets correspond to edges sharing a
vertex in the 2D polygon.
\end{proof}

\begin{lemma}[Infeasibility of $k = 1$]\label{lem:k1-infeasible}
% QC: closure condition for k=1 requires sum of at most 2 same-type
%   normals to vanish, which needs antiparallel normals. Adjacent
%   edges of a convex polygon have normals with angle < pi.
% Downstream: skip k=1 in enumeration. Start at k=2.
For a nondegenerate Lagrangian product, no simple Reeb orbit with
the structure of Lemma~\ref{lem:sigma-structure} has $k = 1$.
\end{lemma}

\begin{proof}
A $k = 1$ orbit has $\sigma = (Q\text{-block}, P\text{-block})$.
The closure condition $\sum_i \beta_i n_{\sigma(i)} / h_{\sigma(i)} = 0$
(Remark~\ref{rem:simple-orbit-data}, applied to
the constraint set~\eqref{eq:constraint-set})
decomposes into independent $q$- and $p$-components.

The $q$-component is $\sum_{i \in Q\text{-block}}
(\beta_i / h_i)\, \hat{n}_i^q = 0$.
A $Q$-block contains one or two $q$-facets.
\begin{itemize}
\item One $q$-facet: $(\beta / h)\, \hat{n}^q = 0$ requires
  $\beta = 0$, contradicting $\beta > 0$.
\item Two adjacent $q$-facets: $(\beta_1 / h_1)\, \hat{n}_1^q
  + (\beta_2 / h_2)\, \hat{n}_2^q = 0$ with $\beta_1, \beta_2 > 0$
  requires $\hat{n}_1^q$ and $\hat{n}_2^q$ to be antiparallel.
  But these are outward normals of adjacent edges of a convex polygon;
  the angle between adjacent outward normals of a convex polygon
  is strictly less than~$\pi$, so they cannot be antiparallel.
\end{itemize}
In both cases, no positive $\beta$ solution exists.
\end{proof}

\label{subsec:bounce-bound}

\begin{definition}[Minkowski billiard trajectory]\label{def:billiard-trajectory}
% QC: standard definition from Gutkin-Tabachnikov (2002).
%   A closed polygon inscribed in K_q with reflection law from K_p.
% Downstream: the billiard trajectory is the q-projection of a
%   generalized Reeb orbit on bdry(K_q x K_p).
Let $K_q, K_p \subset \mathbb{R}^2$ be convex bodies with
$0 \in \operatorname{int}(K_q)$ and $0 \in \operatorname{int}(K_p)$.
A \emph{closed $K_p$-billiard trajectory in $K_q$} is a closed
polygon $\eta = (x_1, \ldots, x_k)$ with vertices
$x_i \in \partial K_q$ (the \emph{bounce points})
and edges $\Delta_i \coloneqq x_{i+1} - x_i$ (indices cyclic).

The \emph{$K_p^\circ$-length} is
\[
  \ell_{K_p^\circ}(\eta) \coloneqq
  \sum_{i=1}^{k} h_{K_p}(\Delta_i),
\]
where $h_{K_p}(v) = \max_{y \in K_p} \langle v, y \rangle$
is the support function of~$K_p$
(Definition~\ref{def:support-function}, restricted to $\mathbb{R}^2$).
\end{definition}

\begin{theorem}[Billiard characterization of EHZ capacity]%
\label{thm:billiard-characterization}
% QC: proved in Artstein-Avidan & Ostrover (2014) for smooth strictly
%   convex bodies, extended by Rudolf (2022) to polytopes.
%   See papers/hko2024/counterexample.tex Section 2, Theorem 2.13.
% Downstream: this connects the Reeb orbit computation (Section 1)
%   to the billiard enumeration below.
\leavevmode\newline
\textup{(Artstein-Avidan \& Ostrover~\cite{AAO2014};
Rudolf~\cite{Rudolf2022}.)}
For a Lagrangian product $K = K_q \times K_p$:
\[
  c_{\mathrm{EHZ}}(K_q \times K_p)
  = \min\bigl\{ \ell_{K_p^\circ}(\eta) :
  \eta \text{ is a closed $K_p$-billiard trajectory in } K_q \bigr\}.
\]
\end{theorem}

\begin{theorem}[Bounce bound]%
\label{thm:bounce-bound}
% QC: Bezdek & Bezdek (2009) Lemma 2.1 + 2.4: any closed convex curve
%   that cannot be translated into int(K_q) can be reduced to at most
%   n+1 = 3 vertices (for K_q in R^2). Rudolf (2022) Theorem 4
%   applies this to billiard trajectories.
%   See papers/hko2024/counterexample.tex line 426.
% Downstream: k <= 3 in the enumeration. Total |sigma| <= 12.
\leavevmode\newline
\textup{(Bezdek \& Bezdek~\cite{BezdekBezdek2009}; Rudolf~\cite{Rudolf2022}.)}

For convex polygons $K_q, K_p \subset \mathbb{R}^2$,
the minimum in Theorem~\ref{thm:billiard-characterization}
is achieved by a billiard trajectory with at most $3$ bounce points.
\end{theorem}


%%% Proof of billiard algorithm completeness

\begin{theorem*}[Theorem~\ref{thm:billiard-completeness}, restated]
Let $K = K_q \times K_p$ be a Lagrangian product.
Then $c_{\mathrm{EHZ}}(K)$ equals the minimum of $A(S, \sigma)$
(equation~\eqref{eq:action-formula})
over all $(S, \sigma)$ with $\sigma$ matching the pattern
$\bigl([Q \mid QQ]\,[P \mid PP]\bigr)^k$
for $k \in \{2, 3\}$.
\end{theorem*}

\begin{proof}
\textbf{Overview.}\quad
Every valid $(S, \sigma)$ from the restricted set is also a valid
input to the algorithm of Section~\ref{subsec:general-case-algorithm}, so
\begin{equation}\label{eq:billiard-lb}
  \min_{\text{restricted}} A(S, \sigma)
  \geq \min_{\text{all}} A(S, \sigma)
  = c_{\mathrm{EHZ}}(K).
\end{equation}
It remains to show the reverse inequality.

\medskip\noindent
\textbf{Step 1: simple minimizer has alternating block structure.}\quad
By Theorem~\ref{thm:simple-minimizer}, $c_{\mathrm{EHZ}}(K)$ is
achieved by a simple Reeb orbit $\gamma^*$ with representation
$(\sigma^*, \tau^*)$, all $\tau_i > 0$.
By Lemma~\ref{lem:sigma-structure}, $\sigma^*$ has the form
$\bigl([Q \mid QQ]\,[P \mid PP]\bigr)^{k^*}$ for some $k^* \geq 2$
(Lemma~\ref{lem:k1-infeasible} excludes $k^* = 1$).
Hence
\begin{equation}\label{eq:billiard-ub1}
  \min_{\text{alternating-block}} A(S, \sigma) \leq A(S^*, \sigma^*)
  = c_{\mathrm{EHZ}}(K).
\end{equation}

\medskip\noindent
\textbf{Step 2: bound $k^* \leq 3$.}\quad
By the billiard characterization
(Theorem~\ref{thm:billiard-characterization}), $c_{\mathrm{EHZ}}(K)$
is also the minimum $K_p^\circ$-length over billiard trajectories
in~$K_q$. By the bounce bound (Theorem~\ref{thm:bounce-bound}),
this minimum is achieved by a trajectory with at most $3$ bounce points.
Each bounce point corresponds to one $q$-block in the simple orbit
(Lemma~\ref{lem:sigma-structure}), so $k^* \leq 3$.

\medskip\noindent
\textbf{Conclusion.}\quad
Combining~\eqref{eq:billiard-lb} and~\eqref{eq:billiard-ub1}
with the restriction $k \in \{2, 3\}$:
\[
  c_{\mathrm{EHZ}}(K)
  \leq \min_{k \in \{2,3\}} A(S, \sigma)
  \leq \min_{\text{alternating-block}} A(S, \sigma)
  \leq c_{\mathrm{EHZ}}(K). \qedhere
\]
\end{proof}

\begin{remark}[Complexity]\label{rem:billiard-complexity}
% Downstream: total iterations for k=2,3 combined.
%   For each k: choose k q-blocks from O(n_q) blocks, k p-blocks
%   from O(n_p) blocks, then enumerate orderings.
%   Dominated by k=3: O(n_q^3 * n_p^3) iterations.
%   Each iteration: solve (|sigma|+5) x (|sigma|+5) system with
%   |sigma| <= 12, so each solve is O(1).
Let $n_q = F_q$ and $n_p = F_p$ denote the edge counts of $K_q$
and $K_p$.
The number of $(S, \sigma)$ to evaluate is
$O(n_q^3 \, n_p^3)$, dominated by $k = 3$:
\begin{itemize}
\item Choose $3$ non-overlapping $q$-blocks: $O(n_q^3)$ choices.
\item Choose $3$ non-overlapping $p$-blocks: $O(n_p^3)$ choices.
\item Orderings per selection: at most $2! \times 3! = 12$
  (fix first $q$-block to remove cyclic symmetry,
  permute remaining $q$-blocks and all $p$-blocks).
\item Within-block orderings: at most $2^6 = 64$
  (each of 6 blocks can have reversed internal order).
\end{itemize}
Each evaluation solves a linear system of size at most
$12 + 5 = 17$, which is $O(1)$.
The total cost is polynomial in $n_q$ and $n_p$.
\end{remark}
